% Standalone methodology draft for https://github.com/njzhx213/qemu-passes
% Repository snapshot consulted: main, as presented by GitHub on 2026-10-04.
% This file deliberately lives outside that repository so it can be downloaded
% and reviewed before being copied into the paper source tree.

\section{Methodology}\label{sec:methodology}

\subsection{Framework overview}\label{sec:framework-overview}

Our framework converts synthesizable peripheral RTL into executable device
models for firmware rehosting.  Its inputs are SystemVerilog RTL, elaboration
parameters, interface bindings, and SoC-level connection metadata.  CIRCT's
\texttt{circt-verilog} frontend elaborates the RTL and produces MLIR containing
LLHD and HW dialect operations.  A sequence of normalization, analysis, and
lowering passes turns that representation into a flat, stateful C model.  A
generated QEMU wrapper exposes the model as an MMIO device and connects its
non-MMIO interfaces to reusable adapters (called \emph{organs} in the
implementation).  SoC generation then instantiates the devices at their
platform addresses and wires interrupts, pins, alerts, and peer buses.

The pipeline has two complementary modeling paths.  The \emph{execution path}
retains the elaborated RTL state and control logic and emits an executable C
transition system.  The \emph{behavioral path} retains the RTL register block
but replaces a selected internal core with a declarative behavior table and a
shared runtime.  The former maximizes reuse of RTL semantics; the latter is a
deliberate abstraction for cores whose full cycle-level implementation is
unnecessary or impractical for the intended firmware experiment.  Both paths
join at the same QEMU-facing interface and are evaluated with the same layered
validation discipline.

\figplaceholder{RTL, parameters, bindings, and SoC metadata entering a CIRCT/
LLHD pipeline.  Split the normalized IR into (i) execution-model lowering and
(ii) register-block plus behavior-table lowering.  Join both at generated C
models, reusable QEMU organs, SoC glue, and firmware validation.  Mark generated,
declarative, shared-runtime, and handwritten artifacts with distinct styles.}
{Architecture of the RTL-derived firmware-rehosting framework.}
{fig:architecture}

We treat generation as a chain of explicit contracts rather than as a single
``RTL-to-QEMU'' claim.  The frontend contract is successful elaboration with the
intended parameters.  The compiler contract is that every retained operation is
either lowered or reported, never silently discarded.  The binding contract is
that every externally relevant port has a declared role.  The integration
contract is that generated instances and their connections match the selected
SoC.  Finally, passing firmware is evidence only for the paths that the firmware
actually exercises.

\subsection{T1: execution semantics translation and behavioral abstraction}

\subsubsection{State representation and scheduling contract}

For a peripheral instance, the emitter produces a C state structure containing
registers, signals, memories represented by the supported aggregate forms, and
flattened input and output leaves.  Let $s_k$ denote clocked state immediately
before model step $k$, $u_k$ the externally supplied inputs, and $c_k$ the
combinational values after settling.  The intended schedule is
\begin{align}
  c_k &= G(s_k,u_k), \\
  s_{k+1} &= F(s_k,c_k,u_k).
\end{align}
The generated \texttt{update\_state()} evaluates combinational logic in
dependency order.  Strongly connected components that cannot be topologically
ordered are evaluated with a bounded local fixed-point procedure.  The generated
\texttt{tick()} computes all next-state values into temporary
\texttt{\_qp\_next\_*} storage before committing them.  This two-phase update is
essential: sequential assignments that are simultaneous in RTL must not acquire
an ordering merely because their C statements are ordered.

This schedule is an executable contract, not a formal equivalence proof.  The
lowering must additionally preserve bit widths, signedness, slices,
concatenations, packed and unpacked aggregate layout, dynamic indexing, instance
specialization, and port connectivity.  Unsupported constructs and unresolved
bridges are therefore surfaced by translation audits rather than converted to
plausible default values.  Wide values are handled as word arrays when they do
not fit the host scalar representation; partial writes to the same word use the
pre-edge or already-staged value required by the source semantics.

\subsubsection{MMIO transactions and bounded settling}

Generated read and write entry points translate a QEMU MMIO operation into a
TL-UL request.  An access asserts the request, evaluates combinational logic,
advances one request clock, deasserts the request, and then invokes the common
settle loop.  The loop repeatedly evaluates and ticks the model until one of
three exits occurs:
\begin{enumerate}
  \item \textbf{fixed point:} a tick changes no sequential state;
  \item \textbf{revisited state:} Brent cycle detection finds a repeated
        fingerprint of the complete generated sequential state; or
  \item \textbf{budget exhaustion:} the access consumes the configured budget
        (65,536 ticks in the evaluated repository snapshot).
\end{enumerate}

The fixed-point exit is normal quiescence.  A repeated state is reported as an
autonomous oscillation for the fixed external inputs: because the comparison
covers the generated sequential state rather than only selected progress bits,
it detects periods longer than two without retaining an unbounded history.
Budget exhaustion is instead a responsiveness boundary.  It is logged, is not
reported as successful device completion, and leaves state available for
continuation on a later access or external event.  This distinction is needed
for engines such as DMA: a long but progressing transfer must not be labeled an
oscillation merely because it exceeds a short callback window.

The fingerprint does not by itself encode changes outside the model.  A reusable
adapter that can synchronously call back into the device must therefore obey the
settle contract.  In particular, the generated busy guard prevents recursive
settling: a re-entrant pin update records the new raw input and the outer loop
observes it on its next iteration.  Tests separately exercise fixed points,
injected cycles, budget continuation, and adapter-driven progress.

\subsubsection{Event clocks and time-dependent state}

Ordinary settling is transaction driven, whereas timers and serial protocols
also depend on events that occur without a CPU MMIO access.  The model retains
edge conditions for externally observable clocks, and adapters inject changes
to clock, pin, or serial input leaves.  Virtual-time pump adapters advance the
free-running timer models while the guest is waiting, allowing timer interrupts
to wake a CPU in \texttt{wfi}.  Clock-ratio state converts a QEMU scheduling
interval into the device-clock advances used by the selected instance.

Pump policy is part of the evaluated configuration rather than an inherent
property of LLHD lowering.  For every time-dependent device we record the pump
source, nominal ratio, batching rule, and stopping condition.  Consequently, the
framework claims functional event ordering and the tested interrupt milestones,
not cycle-accurate correspondence between QEMU virtual time and silicon time.

\subsubsection{Register-block behavioral abstraction}

The alternative path keeps the generated register interface and expresses the
abstracted core as a table interpreted by a shared runtime.  A table row can
apply an immediate register effect, schedule a deferred event, or access a
storage window.  Abstractly, the runtime implements
\begin{equation}
  (q,a,e,t) \longmapsto (q',E'),
\end{equation}
where $q$ is quiescent device state, $a$ an optional bus access, $e$ an external
event, $t$ virtual time, and $E'$ the updated deferred-event set.  It preserves
the software-visible contract selected for the experiment, not every internal
clock or microarchitectural state.

Scope selection is based on structural and peer-aware analysis: whether state is
observable through registers or interrupts, whether a transport must remain
live, whether firmware checks a computed result, and whether the peripheral is
part of a cross-device validation group.  The table and its classification are
declared experimental inputs and are versioned alongside bindings.  Where both
an execution model and a behavioral model exist, trace replay compares their
selected MMIO results and externally visible events.  Independent known-answer
tests remain necessary because two outputs produced from the same compiler input
can share a defect.

\subsection{T2: declarative peripheral and system integration}

\subsubsection{Structural digests, bindings, and reusable organs}

Before QEMU code generation, the pipeline emits a digest of flattened ports,
widths, directions, clocks, resets, and candidate interfaces.  A binding file
assigns semantic roles that cannot reliably be recovered from names alone.  The
generator checks the binding against the digest so that stale port names, width
mismatches, missing members, and ambiguous buses fail early.

Bindings select reusable QEMU organs for character I/O, primary and auxiliary
pin groups, interrupts, alert inputs, timers, SPI, I2C, and bus mastering.  They
may also declare intentional constant inputs, for example environmental or
fault-countermeasure signals not supplied by the selected board.  Every such tie
is auditable configuration, not proof that the tied-off behavior is generally
correct.  We report tied inputs and unconnected outputs per device and distinguish
them from ports exercised by a peer.

For DMA, the \texttt{bus\_master} organ observes the generated model's TL-UL host
requests, performs the corresponding access in QEMU's address space, and returns
the response to the model.  Transfer sequencing, address updates, chunking, and
interrupt generation remain consequences of the RTL engine.  The evaluation
records accepted transfer sizes, byte-mask handling, fixed versus incrementing
addresses, address-space selection, and propagation of bus errors; a negative
test supplies an invalid transaction and requires the firmware-visible error
path to fire.

\subsubsection{Bridging signal- and transaction-level interfaces}

Some interfaces cannot be integrated by directly mapping like-named ports.  An
RTL SPI master exposes clocks, chip select, and individual bits, whereas QEMU's
SSI peer exchanges bytes.  The SPI-master organ therefore uses
snapshot--speculate--replay.  At a byte boundary it snapshots model state,
speculatively advances enough edges to recover the outgoing byte, exchanges that
byte with the SSI peer, restores the snapshot, and replays the same edges with
the returned input bits.  The committed trajectory thus contains the peer's
response while the peer transaction is issued only once.

We validate this bridge at two levels: a generated SPI host communicates with a
QEMU flash using an unmodified OpenTitan DIF driver, and a closed loop routes a
generated host through the pin bridge into a generated SPI device.  The method
records supported SPI mode, bit order, chip-select assumptions, and the number
of speculative and replayed edges per transferred byte.  Snapshot/replay is a
functional granularity bridge; it is not a claim of host-time performance or
analog timing fidelity.

I2C requires a different bridge because SDA and SCL are bidirectional open-drain
wires.  The organ resolves all drivers into a shared line value and feeds the
result back to the model.  The self-loop test uses one generated instance in both
host and target roles and checks START, address, acknowledgement, and captured
data.  Pin groups similarly map named generated outputs and inputs through QEMU
GPIO arrays; auxiliary groups cover non-primary buses such as the pinmux
peripheral side.

\subsubsection{SoC metadata and connected execution}

SoC metadata supplies each instance's RTL top, elaboration parameters, MMIO
range, interrupt source, binding, and peer connections.  Using instance
parameters is necessary because a standalone IP default need not match its
Earl Grey instantiation.  Generation produces device declarations, creation and
mapping code, and connection glue; small board-specific adapters remain
handwritten and are reported separately.

Connected-device groups are first-class validation units.  In the evaluated
snapshot these include: UART receive to a DMA hardware-handshake trigger and DMA
bus-master reads into SRAM; DMA streaming a message into HMAC and receiving the
completion interrupt through the generated PLIC; UART alert injection through
the generated alert handler and PLIC; PWM and pattern-generator outputs routed
through pinmux and a pad loopback into GPIO; and SPI and I2C self-loops.  Such
rings test data and control paths that isolated register tests cannot reach.  A
ring is counted as covered only when its endpoint asserts a content-level result
(for example byte-exact memory or digest data), not merely when all components
were instantiated.

The frozen inventory contains 17 translated peripherals and a regression gate
of 54 targets.  This number denotes the repository's named configurations and
tests, not exhaustive feature coverage.  The evaluation table therefore reports
for every peripheral: frontend/register status, backend organ, generated versus
handwritten integration, gate targets, exercised features, tied interfaces, and
known untested features.  In particular, exported pins without a waveform or
peer-level assertion are not labeled end-to-end validated.

\subsection{T3: layered validation for trustworthy rehosting}

No single test establishes RTL equivalence or platform completeness.  We use
four cumulative layers and preserve the failure artifacts from each layer.

\paragraph{L1: compiler and structural audits.}
The compiler records unsupported operations, dropped or complex expressions,
unresolved inputs, undriven outputs, assignment coverage, and connection
consistency.  Generated C is compiled with warnings enabled, and all generated
instances are checked against their binding and SoC metadata.  These checks can
reach code that firmware never executes, but they cannot show that a lowered
operation has the correct semantics.

\paragraph{L2: register conformance.}
Host harnesses and firmware exercise reset values, read/write permissions,
masks, write-one-to-clear fields, FIFO/window behavior, side effects, and
register locks.  Coverage is computed against an explicit register/field
inventory.  A register is ``touched'' when an access reaches it and
``semantically checked'' only when an independent assertion verifies its value
or side effect; the evaluation reports these measures separately.

\paragraph{L3: functional, known-answer, and differential checks.}
Cryptographic blocks use published independent vectors: SHA-256 through HMAC,
AES-128 ECB and CTR, and SHA3/cSHAKE/KMAC modes.  Multi-operation tests check
state transitions such as counter advancement, mode changes, prefix/key
handling, and interrupt clearing, rather than only a final single operation.
When a reference implementation exposes intermediate rounds, the harness
compares the first diverging state to localize compiler errors.  Differential
trace replay between model forms supplements these known-answer tests but does
not replace them.

\paragraph{L4: real software and connected execution.}
The strongest tests boot the generated devices inside the lowRISC/OpenTitan QEMU
machine and drive them with OpenTitan DIF code or an official smoke test.  They
cover character I/O, external peers, DMA, interrupts, alerts, and the connected
rings described above.  Every target has an explicit oracle: exact bytes,
register values, interrupt observation, a published digest/ciphertext, or a
documented boot milestone.  Runtime stand-ins and source changes are listed per
target.  A boot or passing smoke test proves only its named milestone and is not
reported as complete operating-system or peripheral support.

\paragraph{Negative controls and anomaly attribution.}
For each major mechanism, we introduce a fault expected to invalidate a nearby
oracle: corrupt a generated data-path operation, disconnect an interrupt or
alert, invalidate a bus address, alter a byte mask, or inject an autonomous
oscillation.  A test contributes mechanism evidence only if the corresponding
negative control fails.  Crashes and hangs are classified using the settle exit,
generated trace, QEMU log, and an independent reference; they are not attributed
to guest firmware solely because they occur while running that firmware.

\subsection{Reproducibility protocol and validity boundaries}

Experiments pin the qemu-passes revision, CIRCT revision and build type,
OpenTitan/lowRISC QEMU revision, RTL source revision, bindings, SoC description,
compiler flags, firmware source revision, and each test command.  Generated
artifacts are rebuilt from those inputs rather than treated as ground truth.  A
machine-readable manifest stores these identifiers together with the per-device
model hash and gate outcome.  We archive the complete stdout/stderr log and the
first divergent trace for failures.

The reported results establish that the tested RTL-derived models preserve the
software-visible behavior required by the named tests.  They do not establish
formal equivalence to RTL, cycle accuracy, analog behavior, complete OpenTitan
coverage, or correctness of unexercised modes.  Constant-tied environmental and
security inputs narrow the modeled threat environment.  Frontend workarounds or
semantics-preserving RTL overlays are enumerated and versioned because they are
part of the experimental input.  Finally, shared compiler and generator code can
create correlated errors across devices; independent standards-based vectors,
external QEMU peers, official firmware, and negative controls are used to reduce
rather than eliminate this risk.
